% =====================================================================
%  Professioni 3 — Il narrative designer
%  Fonte: trascrizione della presentazione del corso di narrative design
%  ai Dev Days (AIV3): un incontro online su Zoom di circa un'ora e
%  mezza, il martedì mattina dalle 10 alle 11:30, tenuto da Daniele
%  Fusetto, direttore del corso di narrative design di AIV. Una ventina
%  di partecipanti di tutti i corsi del primo anno; slide condivise; una
%  mezz'ora abbondante di domande, a voce e in chat.
%  Pacchetti richiesti nel preambolo del libro:
%    \usepackage[utf8]{inputenc}, \usepackage{booktabs},
%    \usepackage{tikz} con
%    \usetikzlibrary{arrows.meta,positioning,calc,fit}
%
%  Collocazione nel libro: parte "Alcune Professioni nel Mondo
%  Videoludico", dopo AIV2_game_e_level_design (3.4).
%  Tutte le etichette di questo capitolo hanno il prefisso "p3:".
% =====================================================================

\chapter[Il narrative designer]{Il narrative designer: comunicare
attraverso narrazione e gameplay}
\label{p3:cap:narrative}

La seconda presentazione di corso dei Dev Days (capitolo~\ref{p1:cap:devdays})
è quella del corso di \textbf{narrative design}, tenuta dal suo direttore,
\textbf{Daniele Fusetto}, che si era già presentato il giorno prima (``mi
sono già presentato ieri, quindi sapete più o meno chi sono''). Come tutte le
presentazioni della settimana, non è rivolta solo agli studenti del corso:
``è pensata sia per i ragazzi di narrative, ma soprattutto per chi
collaborerà con loro nei prossimi anni''. Agli studenti di narrative il
relatore fa notare che vedranno ``delle cose che poi rivedremo nel dettaglio
durante il corso'', e che questa è ``l'unica occasione'' in cui lo sentiranno
fare lezione in italiano: il corso è interamente in inglese.

Il relatore spiega perché ci tiene che tutti abbiano una base: negli anni
``ci sono stati parecchi problemi''. I narrative designer sono stati a lungo
visti come ``quelli stronzi, cagacazzo, che facevano le punte alle matite'',
e ``purtroppo a ben donde'': soprattutto all'estero, alcune personalità
hanno alimentato l'immagine del narrative designer come qualcuno che vuole
fare il lavoro di tutti. Negli ultimi cinque anni, dopo il COVID, le cose
sono cambiate, ma restano perplessità e incomprensioni su questa figura.

L'incontro si tiene online, su Zoom, il \textbf{martedì mattina dalle 10
alle 11:30}, con ogni probabilità il \textbf{6 ottobre 2026}\footnote{Per la
deduzione della data si veda il capitolo~\ref{p1:cap:devdays}. Al termine il
relatore fissa la prima lezione di narrative a ``venerdì 9 ottobre''.}. I
partecipanti sono meno del giorno prima (``una ventina''), e il relatore ne
saluta molti per nome. La presentazione è ``molto corta'', per lasciare
spazio alle domande: e infatti la mezz'ora finale è tutta di domande, molte
delle quali, sull'intelligenza artificiale e sullo stato dell'industria,
vanno ben oltre il narrative design.

Avvertenza di lettura: la trascrizione è automatica e storpia moltissimi
nomi propri (``Stephen Dynart'' per Stephen Dinehart, ``Rosberger'' per Ross
Berger, ``Engage Nokia'' per N-Gage, ``manifesto Giacomo'' per manifesto
Agile, ``Siri'' per Ciri, ``Edis 2'' per \emph{Hades II}); dove la
ricostruzione è sicura il capitolo usa la forma corretta, altrimenti lo
segnala in nota. Diverse affermazioni del relatore su fatti storici o di
cronaca dell'industria sono riportate come le dice; dove risultano inesatte,
o non sono verificabili, lo si segnala in nota.

Il percorso del capitolo:
\begin{enumerate}
  \item il corso di narrative design in AIV (sezione~\ref{p3:sec:corso});
  \item le \textbf{definizioni}: chi ha inventato il ruolo, che cosa dice
        Ross Berger, la definizione del relatore
        (sezione~\ref{p3:sec:definizioni});
  \item comunicare a un pubblico preciso: il caso di \emph{Among Us}
        (sezione~\ref{p3:sec:pubblico});
  \item narrative designer e game writer, funzioni e carriera
        (sezione~\ref{p3:sec:carriera});
  \item i \textbf{falsi miti} e il rapporto con il game design
        (sezione~\ref{p3:sec:miti});
  \item lo \textbf{skill set} (sezione~\ref{p3:sec:skill});
  \item il \textbf{mindset} (sezione~\ref{p3:sec:mindset});
  \item le \textbf{cadute} tipiche e il gameplay come re
        (sezione~\ref{p3:sec:pitfalls});
  \item come \textbf{cooperare} con un narrative designer, dove può aiutare,
        le best practice (sezioni da~\ref{p3:sec:cooperare}
        a~\ref{p3:sec:best});
  \item le domande finali (sezione~\ref{p3:sec:domande}).
\end{enumerate}

% ---------------------------------------------------------------------
\section{Il corso di narrative design in AIV}
\label{p3:sec:corso}

Il corso di narrative design esiste in AIV da \textbf{sei anni}. È nato come
``costola'' del corso di game design, quando ci si è accorti che molti
studenti di game design avevano più attenzione e interesse per la narrativa
che per il game design.

\subsection{In inglese, fra arte e tecnica}

Il corso è \textbf{al cento per cento in lingua inglese}, e i suoi studenti
lavoreranno sempre in inglese. Il relatore lo dice ``già subito a tutti gli
altri'': ``abituatevi a lavorare in lingua inglese nel 2026''. Non ne può più
di vedere italiani che rifiutano l'inglese: non ha mai lavorato con un
cliente, anche italiano, che non ne avesse bisogno. Anche i giochi del
secondo anno vanno fatti in inglese, perché \textbf{il mercato in lingua
inglese è cinquanta volte quello in lingua italiana}; forse l'unico mercato
più grande, ``anzi, senza forse'', è quello in lingua spagnola. E per mercato
in una lingua si intende quello di tutti coloro che la parlano, non solo
l'America o la Spagna.

È un corso \textbf{50\% tecnico e 50\% artistico}: gli studenti vedono game
design, vedono l'engine, non fanno programmazione (``ovviamente, perché non
sono dei programmatori''), ma fanno \textbf{arte dentro un contesto tecnico}.
Spesso si scrivono dialoghi non in un \emph{word processor} ma in uno
\textbf{spreadsheet}, in engine, in Unreal: bisogna saper essere creativi in
un ambiente per lo più non creativo.

\subsection{I tre anni}

\begin{description}
  \item[Primo anno.] Il \textbf{branching dialogue}, il dialogo ramificato,
        cioè a più scelte, e le basi della narrazione, oltre ad altri
        moduli.
  \item[Secondo anno.] I \textbf{giochi story driven} e le \textbf{branching
        narratives}: plot ramificati, \textbf{quest design}, \textbf{open
        world design}, ``un po' come succede anche in game design''.
  \item[Terzo anno.] La \textbf{narrativa modulare e procedurale}, ``la crème
        de la crème del nostro lavoro''. Il relatore non è stufo delle storie
        lineari (``una buona storia lineare, ben scritta, mi piace
        tantissimo''), ma la narrativa ``brilla molto di più'' quando c'è una
        qualche forma di non linearità.
\end{description}

Quando si parla di narrativa \textbf{ramificata} (\emph{branching}) si intende
una narrativa in cui l'utente può scegliere dove andare, ``più o meno come
in \emph{Bandersnatch}'' (che però, per il relatore, ``non è un gioco, è
un'esperienza interattiva''): data una domanda, si sceglie se fare A o B, e
questo crea possibilità multiple di cui il narrative designer si deve far
carico, \textbf{anche a livello di budget}: ``un gioco branched costa almeno
due volte e mezzo un gioco non branched''.

\subsection{Le discipline e il portfolio}

Le principali discipline del corso: \textbf{narrative design} e
\textbf{world building}; \textbf{writing}, un modulo-workshop di scrittura in
inglese; \textbf{game design}; \textbf{character design}; \textbf{branching},
soprattutto al secondo anno. Ma anche marketing, \emph{agency}, world
building per i giocattoli (con Simone Penati e Giochi Preziosi), camera in
engine, diritto d'autore al terzo anno, ``tante cose''. ``Quindi non siamo
solo narrativi'': \textbf{narrativo non vuol dire scrittore}.

I narrative designer hanno anche un portfolio particolare: rispetto ai game
designer sono ``una sorta di ibrido'' fra designer e artista. Il loro
portfolio contiene sia \textbf{giochi pubblicati}, ``assolutamente
fondamentali'', sia \textbf{attività non legate ai giochi}: racconti brevi,
romanzi, script, pubblicati o no.

Una slide mostra l'equilibrio, nei tre anni, fra le \textbf{competenze
narrative} e le \textbf{competenze ludiche} (``pratico e di mindset''). Nel
corso di narrative design si fa anche game design, \textbf{per forza}: ``la
differenza fra un bravo narrative designer e uno scrittore che ha la
velleità di dire che è un narrative è che il narrative sa fare game
design''. Non lo fa come un game designer: ha un altro focus, e fa game
design per far sì che il gioco \textbf{comunichi} al meglio. Nelle
competenze narrative c'è molta teoria al primo anno, soprattutto nei primi
sei mesi, e molta pratica al terzo. Nelle competenze ludiche si comincia con
una buona dose di teoria, al secondo anno teoria e pratica si equivalgono, e
al terzo si lavora molto sulla \textbf{mentalità}: magari non si vedono le
teorie più avanzate di game design, ma si impara a mettersi nei panni del
game designer. ``Anzi, mettersi nei panni degli altri, per un narrative
designer, è una delle skill più importanti.''

% ---------------------------------------------------------------------
\section{Che cosa fa il narrative designer: le definizioni}
\label{p3:sec:definizioni}

``Io so che voi siete qua per capire cos'è che fa il narrative designer:
non è che scrive e basta.''

\subsection{Il ruolo più giovane dell'industria}

Il narrative designer è \textbf{la disciplina e il ruolo più giovane della
game industry}, nato fra il 2004 e il 2006. C'è ``una sorta di guerra'' su
chi abbia inventato per primo il nome. La prima offerta di lavoro, del 2006,
fu scritta da \textbf{Stephen Dinehart}\footnote{La trascrizione dice
``Stephen Dynart''. Stephen E. Dinehart IV rivendica di aver coniato il
titolo di \emph{narrative designer} a metà degli anni Duemila. Il giudizio
sulla persona è del relatore.}, ``non una persona di cui ho grande
rispetto'', perché negli anni ha cercato di appropriarsi del termine, che in
realtà non ha creato lui: il termine, dice il relatore, era emerso già nel
2004 alla GDC. Il primo ruolo \textbf{accreditato} di narrative designer è
quello di \textbf{Matthew Toner}, nel 2005, un anno prima, per \emph{Rifts:
Promise of Power}, un gioco per il \textbf{Nokia N-Gage}\footnote{La
trascrizione dice ``Engage Nokia''. \emph{Rifts: Promise of Power} è uscito
nel 2005 per N-Gage; il nome del primo narrative designer accreditato è
quello riferito dal relatore e non è stato verificato.}, ``questo futuristico
cellulare orizzontale fantastico''; ``archeologia del software e
dell'hardware''.

\subsection{La definizione di Dinehart}

La definizione data da Dinehart all'inizio è comunque ``abbastanza ok'':
il narrative designer si focalizza nel fare in modo che il \textbf{gameplay},
l'\textbf{esperienza del giocatore}, il \textbf{livello di gioco} e tutta la
parte di \textbf{storia, narrativa e script} siano \textbf{coerenti},
\emph{exciting} e \emph{compelling}. Quindi non è un ruolo singolo, è un
\textbf{ruolo ibrido}; anzi, molto spesso il narrative designer ha
\textbf{bisogno degli altri} per fare le cose. ``Non so quante volte mi è
capitato di dire: in questo gioco servirebbe per forza una cutscene, solo che
non abbiamo il modello 3D di questo NPC che fa quell'animazione.'' Allora
bisogna chiedere ai modellatori 3D e agli animatori se hanno spazio nella loro
pipeline.

\subsection{La definizione di Ross Berger}

La definizione del relatore deriva da un libro ``importantissimo'', forse il
più tecnico e accademico mai scritto sull'argomento, e uno dei più
dettagliati: \emph{Dramatic Storytelling and Narrative Design} di
\textbf{Ross Berger}\footnote{La trascrizione dice ``Rosberger''. Il libro è
\emph{Dramatic Storytelling and Narrative Design: A Writer's Guide to Video
Games and Transmedia} (2019).}. Berger ha esperienza da narrative designer,
ma è un professore, ``e c'è una grande differenza tra essere professore ed
essere professionista'': il relatore lavora e insegna, e chi solo insegna a
volte non è aggiornato sulle problematiche tecniche dell'industria, su ciò
che si vede giorno per giorno. È la differenza fra ``un generale che sta nel
suo ufficio e un soldato di truppa che va in trincea''; meglio essere
entrambi, ``e ce ne sono''. Berger comunque ``è molto bravo, molto esperto, fa
ricerche molto accurate''.

Secondo Berger fare narrative design è creare \textbf{un percorso per la
storia} (\emph{a pathway for story}) perché venga \textbf{rivelata un pezzo
alla volta}, \textbf{quando certe condizioni sono soddisfatte}, all'interno
dei momenti di gameplay o fra l'uno e l'altro. Una buona definizione, ma
``lacunaria'': dà l'impressione che il narrative designer sia legato per forza
alla storia, e non è così. Basta guardare le \textbf{dieci aree} in cui
Berger stesso dice che il narrative designer lavora: la struttura del plot,
i personaggi, la scrittura ``piana'' (scrivere dialoghi), ma anche il
\textbf{content editing} (l'editing dei contenuti del prodotto), gli
\textbf{story mechanisms} (meccanismi narrativi, non il percorso della
storia), il \textbf{golden path} (che i game designer spiegheranno: non
esattamente il percorso principale, ma il percorso ideale migliore che i
designer hanno pensato per un gioco), la \textbf{IP delivery} (la consegna
della proprietà intellettuale), la \textbf{localizzazione}, la \textbf{voice
over direction}: cose che con il ``percorso'' non c'entrano proprio.

\subsection{La definizione del relatore}

Qual è allora il filo rosso che chiarisce ``una volta per tutte'' che cosa fa
il narrative designer? Per il relatore:
\begin{quote}
\textbf{Narrative design significa comunicare attraverso narrazione,
\emph{fiction} e gameplay}: un messaggio specifico, per un pubblico
specifico, per un obiettivo specifico.
\end{quote}
Il narrative designer è utile \textbf{ogni volta che bisogna comunicare}:
anche un tutorial, anche un cambiamento di regole, anche un cambiamento di
gameplay alla community, anche il fatto che un personaggio è interpretato da
un certo attore. Il relatore cita un caso recente in cui, secondo lui, i
narrative designer non sono stati consultati: in \emph{Marvel Snap} c'è stato
``un po' di trambusto'' per una modalità di draft particolare, che alcuni
giocatori tornati nell'app dopo due anni non avevano mai visto e che nessuno
spiegava, né nel gioco né altrove. ``Ecco, la comunicazione è importante.''

Comunicare attraverso narrazione e gameplay include \textbf{l'invenzione, la
pianificazione, il posizionamento, lo sviluppo di contenuti e di
documentazione} della comunicazione. Per questo il narrative designer è
ibrido, fa sia design sia scrittura: molto spesso le aziende non possono
permettersi di pagare sia un narrative designer sia un game writer, e allora
accorpano le figure.

\subsection{Due libri}

A una domanda in chat (il libro di Berger va bene solo per un narrative
designer o anche per un game designer?), il relatore risponde: ``60\%
narrative designer e 40\% game designer''. E consiglia a entrambe le figure,
``a priori'', un libro scritto da un'esperta di settore: \emph{Narrative
Design for Mobile and Live Games}\footnote{Titolo e autrice (``Valentina
Tamer'') sono riportati come appaiono nella trascrizione e non sono stati
verificati.}. ``E voi direte: Daniele, ma è un libro per i giochi mobile e
live, MMO e quant'altro.'' Sì, ma il titolo ``tradisce'': è un libro sul
narrative design in generale, e agli studenti di narrative il relatore ne dà
le prime cinquanta pagine come dispensa. E poi non si sa dove si andrà a
lavorare: molto spesso ci si ritrova su giochi mobile o online.

% ---------------------------------------------------------------------
\section{Un messaggio, un pubblico, uno scopo}
\label{p3:sec:pubblico}

``Un messaggio specifico per un pubblico specifico per un obiettivo
specifico'' è un'idea che viene dal marketing: \textbf{non si può comunicare
il proprio gioco a tutti}, non può piacere a tutti, ``non potete fare un
gioco per tutti, è inesistente''. Fare un gioco per tutti vuol dire farlo per
otto miliardi di persone, ``e io a malapena riesco ad accordare quattro amici
il sabato sera per andare in birreria''. Bisogna capire \textbf{qual è il
proprio messaggio, il proprio pubblico e il proprio scopo}: solo così si ha
successo.

L'esempio è \emph{Among Us} (``lo conoscete tutti? Mi fate sì con la testa,
vi vedo tipo: oddio, questo ha iniziato a parlare alle 10 del mattino''). Al
lancio non era per niente famoso, era giocato pochissimo: il relatore mostra
il grafico dei giocatori, con il gioco ``poverino'' nel 2019. Poi ``oh oh,
COVID'': gli streamer hanno cominciato a giocarlo in diretta e
improvvisamente il gioco \textbf{ha trovato il suo pubblico}. Senza il COVID,
\emph{Among Us} e la sua azienda ``sarebbero stati morti, letteralmente''.
Invece ancora oggi fa circa \textbf{11.000 giocatori giornalieri}: ``ci sono
giochi per cui i direttori commerciali piangerebbero per avere questi
risultati''\footnote{\emph{Among Us} (Innersloth) è uscito nel 2018; il boom
è dell'estate-autunno 2020. La cifra dei giocatori è quella mostrata dal
relatore.}.

% ---------------------------------------------------------------------
\section{Narrative designer, game writer e carriera}
\label{p3:sec:carriera}

\subsection{Narrative designer e game writer}

Una slide mette a confronto ciò su cui possono lavorare il \textbf{game
writer} e il \textbf{narrative designer}. Il game writer fa \textbf{pura
scrittura}; il narrative designer ha tutta una serie di ambiti su cui il game
writer non lavora: la \textbf{lore} del mondo, la \textbf{story}, i
\textbf{trailer}, la \textbf{bibbia}, il \textbf{plot}. Il game writer
lavora raramente sul plot: scrive, non progetta, e aiuta il narrative
designer nel design.

\subsection{Le funzioni del narrative designer}

Le funzioni del narrative designer, nella slide:
\begin{itemize}
  \item la \textbf{narrative direction}: tenere il tono del gioco, i temi,
        il conflitto centrale, essere \textbf{i campioni della narrativa}.
        (Qui il relatore scherza sulla propria elencazione a punti, ``sono un
        po' Salvini, scusatemi: questo, quello, quell'altro''.) Essere
        campioni della narrativa vuol dire \textbf{difendere il fatto che
        ogni gioco ha una narrativa};
  \item lo \textbf{story design} e il \textbf{world building};
  \item i \textbf{sistemi narrativi} e il narrative design in senso
        stretto;
  \item la \textbf{scrittura} e l'\textbf{editing};
  \item anche la \textbf{direzione degli attori}, perché molto spesso i
        giochi sono interpretati da attori reali.
\end{itemize}

Sugli attori il relatore mostra uno dei primi giochi realizzati in
\textbf{motion capture}, \emph{Final Fantasy VIII}. ``Anzi, il primissimo
motion capture della storia è \emph{Prince of Persia} degli anni Settanta'',
perché le animazioni 2D erano state fatte digitalizzando una persona ripresa
mentre si muoveva, ``girata col Super 8''. ``Cose che si scoprono, ragazzi!
Benvenuti nella tana del Bianconiglio''\footnote{Due precisazioni.
\emph{Prince of Persia} di Jordan Mechner è del 1989, non degli anni
Settanta; e la tecnica usata non era motion capture ma
\emph{rotoscoping}: Mechner riprese il fratello mentre correva e saltava e
ricalcò i fotogrammi. \emph{Final Fantasy VIII} (1999) usò effettivamente
la motion capture per le sue sequenze animate, ma non fu fra i primi giochi a
farlo.}.

\subsection{La gerarchia e la carriera}

La gerarchia classica, ``e sicuramente adesso cambierà con l'IA'', partiva
dal \textbf{junior game writer}, poi \textbf{game writer} e \textbf{senior
game writer}; poi si diventava \textbf{narrative designer}, \textbf{mid},
\textbf{senior}, \textbf{lead} e infine \textbf{narrative director}. Oggi non
è più così: le posizioni sono poche per i narrative designer, ma anche per i
game designer (``tutte le posizioni sono poche, purtroppo''), e quelle da
junior game writer ancora meno, perché le aziende vogliono pagare una
persona per fare più cose: ``se posso permettermi un narrative designer che
può fare anche design, allora il game writer non mi serve'', se non in certe
situazioni. Si vedranno sempre più \textbf{posizioni da freelance} e sempre
meno contratti di tre, sei, nove anni o a tempo indeterminato.

Anche l'esperienza richiesta è cambiata. Un paio di anni fa, per essere
senior bisognava aver lavorato su circa \textbf{sette giochi}; oggi ne bastano
\textbf{tre o cinque}, e il relatore ha visto senior narrative designer con
due giochi all'attivo. Ma va visto quali: due giochini usciti su Nintendo
Switch e giocati da quattro persone, o due MMO di grande successo in sette
anni? ``Prendete queste cose con le pinze.'' Per diventare narrative director
``ci vuole un po'''.

Da qui la domanda che gli fanno spesso ai colloqui: ``se voglio essere la
persona che decide di che cosa parla il gioco, qual è la storia, qual è il
personaggio principale, devo fare game design o narrative?''. ``Ragazzi,
dovete fare gavetta e diventare direttori. Svegliatevi.'' Non è il junior a
decidere su che cosa lavorare: sono il direttore e i capi dell'azienda. Quel
diritto bisogna guadagnarselo.

% ---------------------------------------------------------------------
\section{Falsi miti}
\label{p3:sec:miti}

\paragraph{``Senza storia non c'è bisogno di un narrative.''} Un anno e mezzo
prima, a un convegno, ``in un aperitivo, stavo origliando'', il relatore ha
sentito una scrittrice abbastanza importante in Italia, che stava cominciando
a lavorare come narrative designer (``toh, guarda, non ha mai lavorato come
narrative designer''), dire davanti a una platea di sviluppatori: ``quel gioco
non ha branche, quindi non c'è bisogno di un narrative, non c'è bisogno della
narrativa''.

Il relatore racconta sempre un episodio del 2020, ai tempi di
\textbf{Clubhouse}, il social di chat audio. In una stanza sul narrative
design c'era l'allora lead senior producer di \textbf{King}, quelli di
\emph{Candy Crush}. Ci fu un momento ``awkward, favoloso'': una narrative
designer senior fece un commento piccato, ``volevo spiegare che cos'è il mio
lavoro, perché abbiamo qua Steven, il lead producer di King, che magari non
sa che cos'è un narrative designer''. E lui: ``ma tu conosci King? Perché a
\emph{Candy Crush} io non ho un narrative designer. Ho un team di dodici
narrative designer'' (``poi, con i licenziamenti, non so quanti ne sono
rimasti''). Il punto: \textbf{non è un ruolo di plot, è un ruolo di
comunicazione}. Se si vuole che il gioco faccia capire al giocatore di andare
da A a B, serve un coordinamento fra narrativa e level design; ``e non vuoi
fare la figura di \emph{Wolverine}, con la traccia che ti guida
automaticamente?''\footnote{Nella trascrizione: ``non vuoi fare la figura di
merda di Wolverine con il trail del set che ti guida automaticamente''. Il
riferimento, con ogni probabilità, è a un sistema di guida del giocatore in
\emph{Marvel's Wolverine} di Insomniac che il relatore critica; i dettagli
non sono ricostruibili dalla trascrizione.} Allora serve un narrative
designer.

\paragraph{``Un level designer e un narrative designer si occupano solo di
dialoghi e plot.''} No: c'è molto che va oltre la parola. C'è tanta narrazione
anche nei film senza parole: \emph{A Quiet Place} è un film ``con trenta
pagine di script, perché è tutto muto''\footnote{Il numero di pagine è del
relatore. Lo script del film ha pochissimi dialoghi, e molte pagine di sole
indicazioni d'azione.}. Al videogioco di \emph{A Quiet Place} ha lavorato,
con altri, Giacomo Masi, amico del relatore e narrative designer (``forse ha
lavorato più a un altro progetto, comunque avevano dei team'') per fare il
narrative design di quel gioco.

\paragraph{``Siamo solo una sottobranca del game design.''} Non proprio:
narrative e game design ragionano in modi molto simili, ma non identici. Il
relatore precisa di essere ``della categoria per cui un gioco è prima di tutto
gameplay'', e i narrative designer devono ``mettersi il cuore in pace''.

\subsection{Game design e narrative design: finalità diverse}

C'è spesso il \emph{clash} fra game design e narrative design: ``possono fare
la stessa cosa? Sono la stessa cosa?''. No. \textbf{Condividono strumenti e
pipeline, ma hanno finalità diverse}:
\begin{itemize}
  \item il \textbf{game design} cura la \textbf{funzionalità} e le
        \textbf{meccaniche} per creare una \textbf{risposta emotiva};
  \item il \textbf{narrative design} guarda al \textbf{significato} e al
        \textbf{messaggio} di quella comunicazione emotiva.
\end{itemize}
Se si vuole spaventare il giocatore, ma il gioco è un gioco comico per
bambini e ci si è messo del sangue, ``Houston, abbiamo un problema''.

% ---------------------------------------------------------------------
\section{Lo skill set del narrative designer}
\label{p3:sec:skill}

``Fin qua tutto bene? Mi state seguendo? Ci sono parole che non avete
capito?'' Poi lo \textbf{skill set}: le competenze che un narrative designer
deve avere. ``E qua c'è il problema: non è poca roba.''

Andrea Morini, il direttore della didattica, gli ha fatto notare di recente
che il suo corso ha più lezioni degli altri: ogni anno i narrative fanno
dalle quattro alle cinque lezioni in più. ``Sì, raga, per far entrare nel
corso tutto quello che serve a un narrative sto facendo veramente fatica.''
Un narrative designer deve sapere:
\begin{itemize}
  \item \textbf{game design};
  \item la \textbf{scrittura tradizionale};
  \item \textbf{scrivere ``in macchina''}, cioè in engine;
  \item l'\textbf{emotional design};
  \item gli strumenti di studio del \textbf{flusso informativo};
  \item l'\textbf{architettura dell'informazione} di UX e UI.
\end{itemize}
I narrative designer ``rubano'' a tutte le altre discipline le tecniche che
servono per comunicare al meglio.

E serve anche una fortissima \textbf{capacità di comunicazione}, verso il team
e dentro il team: ``non sapete quante volte mi sono trovato a fare il maestrino
dell'asilo, ma è un po' il nostro ruolo''. ``Non siamo i capi'', dice il
relatore\footnote{La trascrizione dice ``non siamo i char noi'': dal contesto
e da un passaggio successivo (``soprattutto i char, che vogliono fare il
lavoro di altri''), il senso sembra essere ``non siamo i capi''; la parola non
è ricostruibile con certezza.}: siamo comunicatori. Fra le soft skill,
il narrative designer \textbf{traduce ciò che uno dice all'altro}: molto
spesso un game designer, un 3D artist e un programmatore dicono qualcosa che
fa arrabbiare un'altra persona, e il narrative designer è quello che dice
``ma loro due stanno dicendo la stessa cosa in un altro modo''. Ci sono veri
paradossi di questo tipo; anche i game designer lo impareranno, e a lezione
faranno il \textbf{Production Dictionary}, ``molto, molto utile da questo
punto di vista''.

Questo spiega anche un episodio del lunedì (sezione~\ref{p1:sec:metodo}):
Sara aveva scritto ``incontri'' e il relatore l'aveva corretta, ``sono
corsi''. ``Io non faccio i puntini sulle i perché mi piace farlo'', ma perché
ogni volta che c'è un'\textbf{incoerenza nel team}, se non la si risolve
subito, ``a valle avrete solo giochi come \emph{Concord}, e io non credo che
voi vogliate fare giochi alla \emph{Concord}''.

In sintesi: \textbf{coerenza emotiva attraverso il sistema di gioco}. Ogni
forma di scrittura diventa strumento comunicativo, \textbf{tutto il gioco
parla}. ``Metà narrativa, metà design, metà arte, metà progettazione'', ma
sempre con \textbf{metodo}, dove \textbf{metodologia} significa \textbf{la
scelta del metodo giusto} per sviluppare un prodotto.

% ---------------------------------------------------------------------
\section{Il mindset del narrative designer}
\label{p3:sec:mindset}

Come deve pensare un narrative designer?

\subsection{Il design thinking}

C'è anzitutto una base che il relatore consiglia a tutti, game designer,
programmatori e 3D artist compresi: il \textbf{design thinking}. È una teoria
di sviluppo di prodotto degli anni Settanta, poi cambiata e migliorata,
``molto interessante perché oggigiorno tutti quelli che fanno IA se la sono
dimenticata'', ed era la metodologia migliore per sviluppare un prodotto,
perché il primo passo era \textbf{empatizzare con il pubblico}: capire i
problemi dei giocatori, che cosa vogliono, come vogliono emozionarsi. Chi
compra uno sparatutto in prima persona di guerra vuole qualcosa di molto
diverso da \emph{Super Mario Bros.\ Wonder}, e bisogna darglielo.

Il design thinking ha cinque passi: \textbf{empatizzare, definire, ideare,
prototipare, testare}, che si ripetono in cerchio
(figura~\ref{p3:fig:design}). È, dice il relatore, ``la base addirittura del
manifesto Agile''; chi non lo conosce lo cerchi online e lo legga, esiste
anche in italiano\footnote{La trascrizione dice ``manifesto Giacomo''.
Il \emph{Manifesto for Agile Software Development} è del 2001.}.

\begin{figure}[htbp]
\centering
\begin{tikzpicture}[
  passo/.style={draw, rounded corners, minimum width=2.3cm, minimum height=0.9cm,
                font=\small, align=center},
  >=Stealth]
  \node[passo] (E) at (90:2.3) {Empatizzare};
  \node[passo] (D) at (18:2.6) {Definire};
  \node[passo] (I) at (-54:2.4) {Ideare};
  \node[passo] (P) at (-126:2.4) {Prototipare};
  \node[passo] (T) at (162:2.6) {Testare};
  \draw[->, thick] (E) to[bend left=15] (D);
  \draw[->, thick] (D) to[bend left=15] (I);
  \draw[->, thick] (I) to[bend left=15] (P);
  \draw[->, thick] (P) to[bend left=15] (T);
  \draw[->, thick] (T) to[bend left=15] (E);
  \node[font=\scriptsize, align=center] at (0,0) {mettersi\\ nella testa\\ dei giocatori};
\end{tikzpicture}
\caption{I cinque passi del design thinking, ripetuti in cerchio, come li
presenta il relatore. Ricostruzione schematica.}
\label{p3:fig:design}
\end{figure}

Il primo punto è \textbf{mettersi nella testa dei giocatori}. Il relatore lo
mostra con un \emph{pitch} volutamente pessimo (``sto facendo un pitch di
merda apposta''): il ``direttore mega galattico'' annuncia che il prossimo
gioco, per i prossimi tre anni, sarà uno story driven in prima persona, con
elementi multiplayer, in cui si è un astronauta umano, un militare, disperso
nella galassia di Andromeda piena di alieni, che deve tornare a casa; si
scopre che tutta la galassia è in guerra e che la Terra non lo è perché si
trova dentro un buco nero che la protegge; ma c'è un alieno cattivo che ha
capito come invaderla, perché la Terra è importante per qualche ragione.
Davanti a una richiesta così bisogna \textbf{empatizzare} (che cosa
vorranno i giocatori?), \textbf{definire} i problemi e il gioco al meglio,
\textbf{ideare} soluzioni, \textbf{prototipare} e \textbf{testare}, e
continuare in cerchio.

\subsection{Il pensiero organico}

Serve poi ciò che il relatore chiama \textbf{pensiero organico}, o
\textbf{pensiero concreto}. ``Organico'' viene da un doppiatore, Oreste
Baldini, che gli studenti di narrative conosceranno al secondo e al terzo
anno, e dalla sua idea di \textbf{scrittura organica}: quando si scrive non
si deve pensare alla filosofia, ma al fatto che \textbf{ciò che si scrive, a
un certo punto, un 3D artist lo deve fare, un level designer lo deve fare}.
Storie del tipo ``il mio personaggio affronterà grandi dilemmi interiori,
supererà il suo disturbo post-traumatico e alla fine vincerà i suoi demoni
interiori'': ``sì, ma che gioco è? \emph{God of War} o \emph{The Last of
Us}? Non lo so''. Bisogna pensare in termini di \textbf{azioni e scelte},
\textbf{spazi e tempi}, \textbf{esteriorità e interiorità}: non solo ciò che è
dentro il personaggio, ma anche ciò che è fuori. Nell'ultimo film di
\emph{Resident Evil}, dice il relatore (``l'ho detto forse anche lunedì''),
c'è un grande lavoro sull'azione: tutto il film è portato avanti dalle azioni
del protagonista.

\subsection{Pensare come un regista, o come un dungeon master}

Secondo: \textbf{pensare come un direttore, un regista, o un dungeon master},
con un \emph{quid}, che il relatore dice anche ai game designer: si imparerà
dal gioco di ruolo da tavolo, da D\&D, a fare il dungeon master, ma
\textbf{non si è dungeon master}, si è professionisti. Il dungeon master è un
lavoro hobbistico, un giocatore; o comunque chi lo fa per lavoro fa il dungeon
master, non il narrative o il game designer.

Che cosa bisogna fare?
\begin{itemize}
  \item \textbf{Prevedere percorsi imprevedibili.} Chi ha fatto il dungeon
        master una volta lo sa: si scrive una bellissima storia nei minimi
        dettagli, e nei primi dieci minuti di gioco i giocatori ``ve la
        mandano all'aria'' (``è automatico, è classico''; ``motivo per cui
        non scrivo più avventure dal 2006'').
  \item \textbf{Guidare senza forzare.} ``Potete andare da quel lato, dove
        c'è il tempio della divinità, oppure verso quella strada fra gli
        alberi dove a un certo punto c'è una luce, ma non sapete bene che cosa
        c'è.'' Che cosa scelgono i giocatori? La luce, perché il designer ha
        messo loro qualcosa in testa che li fa andare verso la luce, ``come
        delle farfalle''. Si chiamano \emph{nudge}\footnote{La trascrizione
        dice ``si chiamano Naggia'': con ogni probabilità \emph{nudge}, la
        ``spinta gentile'' che orienta una scelta senza imporla.}, e si
        studieranno al secondo anno.
  \item \textbf{Rendere le regole metafore narrative}: ``la cosa più
        interessante del nostro lavoro''. Ciò che fa più arrabbiare il
        relatore, nei giochi narrative driven ma non solo, è una meccanica
        che va nella direzione opposta a ciò che il gioco richiede a livello
        narrativo: ``è abominevole''.
  \item Ricordare che \textbf{i giocatori sono i protagonisti}. In
        \emph{Baldur's Gate 3} il personaggio che si crea non è davvero il
        protagonista: è dentro un party in cui i protagonisti sono gli altri,
        ma è lui a decidere che cosa fare. Il motivo per cui c'è un
        protagonista creato dal giocatore più una banda di personaggi è che
        lo sviluppatore non sa che faccia, che classe, che caratteristiche
        avrà il personaggio creato, e gestire tutte le possibilità mettendolo
        come personaggio principale sarebbe ``atroce''; allora si costruiscono
        personaggi di \textbf{esperienza vicaria}, che si vedranno nel corso.
\end{itemize}

\subsection{Gameplay: emozione più significato}

\textbf{Gameplay significa emozioni più significato.} Se ci sono solo
emozioni, ci si emoziona, ma si gioca ``per niente''; se c'è solo significato,
si capisce che cos'è il gioco, ma non c'è emozione. Lavorare sui giochi,
nota il relatore, ``sembra di essere degli psicologi'': si lavora più sulle
emozioni che su altro. \textbf{Fare giochi significa provocare emozioni negli
altri, a comando}, quando si vuole, perché le esperienze ludiche sono
proprio questo.

Se il giocatore deve sentirsi forte nei panni di un personaggio, il
personaggio deve sembrare forte. E allora si lavora sul microsecondo
dell'animazione; sul \textbf{feedback aptico}, quello che si sente nel
controller quando si colpisce; sul ritardo di certe animazioni e di certe
\emph{cinematic}; su meccanismi del gameplay che il giocatore non nota e non
conosce. Come il segreto di \emph{Pac-Man}: il relatore non ricorda mai se
valga quando si prende o quando non si prende il pallino grosso; ``fin tanto
che non prendi il pallino grosso, ogni volta che mangi un pallino'' Pac-Man
\textbf{rallenta}, ``tu non te ne accorgi'', ed è fatto apposta perché i
fantasmi possano raggiungerlo\footnote{Il relatore parla di ``un secondo''.
In realtà Pac-Man si ferma per un solo fotogramma (circa un sessantesimo di
secondo) per ogni pallino mangiato, e per tre fotogrammi per ogni pallino
grande; ne risulta che nei corridoi pieni di pallini si muove un po' più
lentamente dei fantasmi.}. Ci sono tante piccole ``chicche'' di game design
che non si notano: è per questo che una delle primissime cose da fare quando
si fanno giochi è \textbf{dissezionarli}.

\subsection{Player fantasy e dissonanza ludonarrativa}

Poi bisogna progettare la \textbf{player fantasy}, ``importantissima'', e
bilanciarla con la \textbf{dissonanza ludonarrativa}: ciò per cui a volte è
impossibile che narrativa e gameplay vadano di pari passo. C'è una sorta di
\textbf{sospensione dell'incredulità}. L'esempio canonico: in
\emph{Resident Evil}, ``anche in \emph{Resident Evil 9}'', il salvataggio
avviene alla \textbf{macchina da scrivere} (il relatore pensava fosse il
telefono, e viene corretto: ``ecco, quella''). Una macchina da scrivere non
salva la vita nel mondo reale, eppure il giocatore non se ne accorge, ``dice
che figata''. Ci sono \textbf{dissonanze peggiori, da correggere}, e
dissonanze \textbf{da accettare}, perché altrimenti il gioco non si può
fare.

\subsection{Emotional delivery}

L'ultimo punto è l'\textbf{emotional delivery}: assicurarsi, insieme al game
designer, che \textbf{l'emozione arrivi al giocatore}. Che l'\textbf{intenzione
drammatica}, cioè che cosa davvero vuole significare ciò che si sta facendo e
dicendo, arrivi e venga capita. ``Solo il testo esiste'', cioè \textbf{esiste
solo il gioco finale}: non c'è nessuno accanto al giocatore a spiegargli, e
non si può pensare che vada a vedersi i tutorial; ``meno il giocatore usa i
tutorial, meglio è''. Bisogna pensare a come il giocatore cambia durante la
partita e a come unire il \textbf{tema} al \textbf{design}. Nel gioco che sta
giocando in questo periodo\footnote{La trascrizione dice ``sto giocando a
contro il WrestleMania'': il titolo non è ricostruibile con certezza;
dal resto dell'incontro, potrebbe trattarsi di \emph{Control Resonant}.}
ci sono dialoghi che non dicono direttamente che cosa sta succedendo, ma lo
fanno capire, e preparano all'impatto: si chiama \textbf{foreshadowing}, ed è
intenzione drammatica. La domanda da farsi è ``il giocatore lo capirà?'', non
perché il giocatore sia stupido, ma perché magari non lo capisce \textbf{perché
siamo noi a descriverlo male}.

% ---------------------------------------------------------------------
\section{Le cadute tipiche}
\label{p3:sec:pitfalls}

Le \emph{narrative pitfalls}, ``cadute di stile'' tipiche dell'industria.

\subsection{``La storia poi si adeguerà''}

``Noi creiamo il gioco, la storia poi si adeguerà al resto.'' Il relatore cita
\textbf{Richard Dansky} (``seguitelo su LinkedIn''): \textbf{la storia è la
parte più adattabile di un gioco, finché non lo è più}, perché non ci si può
svegliare all'ultimo per cambiare una cinematica. Nei tripla A le cinematiche
sono fatte in un engine a parte, con una qualità d'immagine superiore, e hanno
\textbf{una pipeline parallela} a quella del gioco: si fissa uno script, lo si
consegna al team delle cinematiche, e le cinematiche cambiano raramente; se
cambiano, è perché quel team dice che una cosa non si può fare tecnicamente.

Quando un direttore (``non succede molto spesso con un lead narrative, ma
con i game director e i creative director succede'') dice ``poi la storia
l'adeguate'', la prima domanda da fare è: \textbf{quella parte è dentro
cinematiche o cutscene? Ha dialoghi doppiati?} Se sì, bisogna richiamare gli
attori a doppiare, e sono costi: la giornata, il teatro di doppiaggio, il
direttore del doppiaggio. ``Uno non pensa che costino, i giochi, ma costa ogni
singola parola.''

È vero, da un certo punto di vista la storia è più adattabile della
programmazione, del game design, del 3D: se il problema tecnico viene
segnalato per tempo, si può fare ``la mossa \textbf{Jessica Fletcher}''.
Negli anni Novanta \emph{La signora in giallo} fece un episodio
``fintissimo'' sulla realtà virtuale, con un problema nel branching: un NPC va
in bug, e la soluzione di Jessica Fletcher è togliere dialoghi e battute a
quel personaggio e rifare un po' la storia, e i designer dicono che in effetti
potrebbe funzionare\footnote{Il relatore racconta l'episodio a memoria;
titolo e dettagli non sono riportati.}. ``Fare una Jessica Fletcher è molto
figo. Si può fare, ma non sempre.'' Con i dialoghi si possono aggiustare
cose, \textbf{ma più si cerca di aggiustare con la storia, più la storia
sembrerà \emph{cringe}}: ci si ritrova con dialoghi del tipo ``ah sì, il
cacciavite dimensionale che abbiamo trovato nella quest numero X, quando hai
fatto queste cose\ldots'', e il giocatore pensa ``l'ho fatto, non c'è bisogno
che me lo rispieghi''.

\subsection{``Io gioco e vedo film, so bene come si narra''}

``Se avessi un centesimo per ogni direttore che mi ha detto questa frase,
sarei ricco sfondato.'' Non è vero, altrimenti chiunque fa il caffè saprebbe
costruire una macchina del caffè, e chiunque guida saprebbe costruire
un'automobile. Giocare e vedere film \textbf{aiuta} (e nei suoi corsi il
relatore chiederà di vedere e sperimentare molti giochi), ma non basta: per
questo si studia in un'accademia. Detto questo, ``offendere non aiuta''. C'è
qualcosa di tecnico nel lavoro del narrative designer, cose studiate che gli
altri non hanno studiato: ``\textbf{fidatevi del narrative designer se vi dice
che qualcosa non torna}''. Non bisogna partire dal presupposto che ciò che
dicono gli altri non sia vero, e vale anche per 3D artist e programmatori. Il
relatore lavora molto bene con gli artisti, perché ha cominciato la carriera
così: ``fidatevi di un artista, se vi dice che non si può fare''.

\subsection{``La storia è la prima cosa''}

``La storia è la prima cosa, ciò da cui parte tutto.'' Qui il relatore deve
dire di no: \textbf{se si fanno giochi, il gameplay è la prima cosa}, anche in
un narrative driven in prima persona a narrativa lineare. ``Se non c'è
gameplay, non c'è gioco.'' \textbf{Il gameplay è il re}: tutti, 3D artist,
programmatori, animatori, QA, sono lì per fare il gameplay. \emph{Wolverine},
dice, ha avuto ``un marketing della Madonna, ma non c'è gameplay, il gioco è
vuoto, non c'è sensazione''.

\textbf{Ma la storia è l'assassino del re}: è ``il bassista del party''. Quando
c'è si pensa ``vabbè, non è importante''; quando non c'è, si sente
eccome. Una slide mostra giochi in cui la narrativa è a volte importantissima,
a volte no, o in cui non si penserebbe che ci sia una narrativa dietro; e
giochi come \emph{Directive 8020}, quasi solo narrativa con pochissimo
gameplay: ma c'è un minimo di gameplay, c'è \emph{agency}. Che sia stato meno
di quanto i giocatori si aspettassero, e abbia raccolto lamentele, è un altro
discorso; non è \emph{Bandersnatch}, dove il gameplay non c'è proprio. Ci sono
giochi con molta \emph{agency} e giochi con poca: dipende dal contesto e dal
pubblico.

% ---------------------------------------------------------------------
\section{Cooperare con un narrative designer}
\label{p3:sec:cooperare}

La parte successiva, \emph{narrative cooperation}, spiega come cooperare con
un narrative designer, ``come non offenderlo e come non essere offesi''. ``Lo
capite come faccio lezione io: non le mando a dire.'' Il relatore è
freelance, non è assunto da un'azienda: se deve contraddire qualcuno lo fa,
``non in maniera diretta'', ma dicendo le cose come stanno, non ``hai ragione
tu'' perché è pagato dal direttore. ``Se il gioco è una merda, è una merda.''
Tre luoghi comuni da sfatare.

\paragraph{``Il narrative designer mi dice sempre come fare il mio lavoro.''}
Non è così. Esempio: si vuole una cinematica prima di un boss per far capire
che il boss è un NPC incontrato in precedenza, poi trasformato. Allora
\textbf{bisogna seminarlo nel gioco}: non basta mostrare l'NPC una volta e poi
mostrarlo trasformato, perché non ci sarà nessun \emph{pathos}. Succede anche
in molti film recenti: ``vedi un personaggio una volta, lo uccidono, e tu sei
lì: oh no, è morto John, quello che ho visto in una scena a caso''. Il
narrative designer dirà: serve almeno una quest con quell'NPC, che ne
costruisca l'umanità e il rapporto con il giocatore. Se gli si risponde ``non
dirmi come fare il mio lavoro'', lui non sta dicendo come farlo: sta dicendo
\textbf{che cosa serve}. Se non si vuole fare una quest, se ne parla e si
trova qualcos'altro.

\paragraph{``Il narrative è un maniaco del controllo, un precisino.''} ``Raga,
la narrazione è precisione'', come il game design. \emph{Baldur's Gate 3} non
è venuto bene perché qualcuno si è messo lì ``con quattro stutti
cadenti\footnote{Così nella trascrizione; il senso è ``alla buona''.} a dire
boh, facciamolo'': hanno avuto precisione. \emph{The Witcher 3} (il relatore
cita la nuova versione rimasterizzata, uscita con una patch gratuita) non ha
avuto successo lasciando tutto al caso: ci hanno lavorato attentamente. A
volte serve spaccare il capello in quattro. Certo, serve anche un game
designer o un programmatore che controbilanci e dica ``vabbè, non spaccare il
capello in quattro adesso'': è un lavoro di \textbf{bilanciamento del team}.
``Però qualcuno che ti spacca il capello in quattro, ogni tanto, serve.''

\paragraph{``Il narrative dà indicazioni, ma poi possiamo fare noi.''} Spesso
si sottovaluta come il pubblico interpreterà le cose, e qui il relatore
``deve diventare cattivo'': se l'anno prossimo, nei Dev Projects, vedrà
dialoghi messi in engine da game designer o programmatori, non rivisti dai
narrative e all'insaputa dei narrative, ``io vi spezzo le braccia''. Un
narrative non si permetterebbe mai di scrivere codice senza chiedere un
secondo parere, perché non sa farlo; allo stesso modo nessuno si permetta di
inserire dialoghi al posto di un narrative, o di decidere cose di trama, senza
chiedere. Altrimenti poi i game designer vanno da lui a piangere: ``Daniele, la
controroom\footnote{Dal contesto, il gruppo di revisione dei Dev Projects
(si veda il capitolo~\ref{p1:cap:devdays}).} ci ha detto che la trama è una
merda''. ``Certo: l'hai fatta tu, non il narrative, o comunque non hai chiesto
al narrative.'' Bisogna \textbf{usare le proprie persone}: se si ha un
programmatore esperto, gli si chiede; se si ha un 3D artist esperto, gli si
chiede. Non dare nulla per scontato.

% ---------------------------------------------------------------------
\section{Dove può aiutare un narrative designer}
\label{p3:sec:aiuto}

\begin{description}
  \item[Nel pitch del progetto.] Al primo e al secondo anno si faranno molte
        slide, molti \emph{pitch deck}, molte presentazioni. Il narrative è
        esperto di comunicazione: bisogna fare un controllo con lui,
        \textbf{fare il pitch deck insieme}. Non farlo fare solo al narrative
        (``tu sei narrative, falle tu'') e neppure solo a chi narrative non
        è: due teste ragionano meglio di una.
  \item[Dove bisogna colmare una lacuna e generare coerenza.] ``Che cosa
        disegno sullo scudo di un guerriero elfico? Com'è la moda su Proxima
        Centauri nel 2355? E come rendo coerente tutto questo con il resto del
        gioco?''
  \item[Per visualizzare insieme una soluzione narrativa.]
        \textbf{Visualizzare} non significa immaginare con la mente, ma
        \textbf{mappare e pianificare i contenuti di gioco}. Un problema
        narrativo è un problema \textbf{comunicativo}. Esempio: a un certo
        punto il giocatore dovrebbe sapere come usare l'ascia del
        protagonista, ma i playtest mostrano che non la usa mai e la ritiene
        debole, anche se i test di bilanciamento dicono che è potentissima.
        Come fargli capire che l'ascia è importante? Inondarlo di messaggi a
        schermo ``usa l'ascia, usa l'ascia'' è ``il modo migliore per fare
        arrabbiare il giocatore''. Meglio chiedere una soluzione narrativa.
  \item[Per rendere concrete le idee.] Non basta essere creativi e avere
        idee: ``le idee le hanno tutti, la capacità di realizzarle ce l'hanno
        solo i migliori''.
\end{description}

% ---------------------------------------------------------------------
\section{Best practice}
\label{p3:sec:best}

Le best practice per lavorare con un narrative designer, ``ma in generale, le
prime due soprattutto valgono per tutti''.
\begin{enumerate}
  \item \textbf{La visione del gioco è al di sopra di tutti.} A un certo
        punto si farà un documento chiamato \textbf{vision} (o in un altro
        modo): un documento, all'inizio dello sviluppo, in cui ci si accorda
        sul gioco da fare. Non è dettagliatissimo; descrive con parole proprie
        \textbf{dove si vuole arrivare}, l'obiettivo del lavoro, per evitare
        di andare fuori strada. Bisogna lasciare spazio per infondere nel
        gioco la propria arte, e quindi anche ``andare un po' \emph{sideways}'':
        magari non si arriverà esattamente al cento per cento di ciò che si
        voleva, ma molto vicino. Ma bisogna anche ricordarsi a vicenda di
        essere \textbf{le barriere} del proprio lavoro. Se si è deciso di
        fare uno sparatutto in prima persona e arriva un programmatore che
        dice ``ho fatto un bellissimo sistema per la terza persona''; ``ma
        noi siamo un first person shooter''; ``vabbè, a me piace di più la
        terza persona'': ``non è così che si fa un gioco''. Se ci si è
        accordati per fare una cosa, si fa quella.
  \item \textbf{Data beats dogma}: \textbf{i dati battono il dogma}. ``Se
        volete avere ragione, cambiate lavoro.'' Nella game industry nessuno
        vuole avere ragione, e chi vuole avere ragione di solito lavora male.
        Se si vuole \emph{dimostrare} di avere ragione, che è diverso, bene:
        bisogna portare le prove, per due motivi. Se il capo o il collega è
        intelligente, davanti ai dati dirà ``pensavo di no, ma dati alla mano
        hai ragione''; e infatti una delle competenze più importanti è
        \textbf{la capacità e l'umiltà di cambiare idea davanti ai dati}. Se
        invece non si lascia convincere (``una mossa trumpiana: 2.500
        playtester hanno detto che quel livello è noioso, io non sono
        d'accordo''), almeno lo si è detto, e nessuno potrà accusare di non
        averlo fatto.
  \item \textbf{Una storia è la promessa di una conversazione.} ``La più
        complessa da capire'', che si approfondirà nel corso. Ogni volta che
        si fa un gioco o una storia si promette al pubblico una
        conversazione. Anche in un film: se all'inizio si mostra un oggetto,
        si dà per scontato che verrà usato, e si apre una conversazione
        implicita con il pubblico, che comincia a chiedersi dove verrà usato.
        Come la mitragliatrice di \emph{Breaking Bad}: Vince Gilligan, che ha
        scritto la serie, non sapeva nemmeno come usarla; l'ha messa ``nel
        pilot'' e poi pensava che il pubblico se ne sarebbe dimenticato, ma
        gli sceneggiatori della writers' room gli hanno detto che il pubblico
        se ne sarebbe ricordato; e infatti poi ne hanno fatto il
        \emph{payoff}\footnote{La mitragliatrice M60 compare non nel pilot ma
        nel \emph{flash-forward} iniziale della quinta stagione (2012), e
        viene usata nell'episodio finale della serie. Sull'aneddoto della
        writers' room il capitolo riporta la versione del relatore.}.
  \item \textbf{Tutto si può risolvere dialogando.} È un altro ``assioma
        fondamentale'' di AIV, su cui quest'anno il relatore vuole essere
        ``ancora più perentorio'', perché l'anno scorso è stata dura.
        \textbf{Tenersi le cose all'oscuro rende il lavoro più difficile}, e
        tenere le cose all'oscuro vuol dire anche fare \textbf{gossip}, andare
        da un compagno a dire ``guarda, questa persona ha detto questa cosa'':
        ``non è una conversazione, quella''. Se si hanno problemi, si chiede;
        se si hanno dubbi, si domanda; se ci sono questioni da risolvere, si
        risolvono. Se ci si è arrabbiati con un compagno, si ha un anno per
        farsela passare, perché l'anno dopo si farà un gioco con lui. Se una
        persona non piace, bisogna capire durante quest'anno come
        collaborarci comunque, perché nel mondo del lavoro non si possono
        cambiare i colleghi perché stanno antipatici. ``Mattia non mi piace''
        (``nome a caso, non ce l'ho con Mattia''), ``però hai firmato un
        contratto di tre anni, ti pagano anche bene, e sei pure full remote'':
        si trovano dei modi per collaborare. \emph{Savoir-faire}, soprattutto.
        ``Fine della paternale.''
\end{enumerate}

% ---------------------------------------------------------------------
\section{Le domande finali}
\label{p3:sec:domande}

``Domande? Avete trenta minuti per le domande.''

\subsection{A che cosa serve un narrative in un gioco da tavolo?}

In un gioco da tavolo, in che modo può aiutare un narrative designer? In vari
modi. Se il gioco è \textbf{astratto}, probabilmente il narrative designer
lavorerà più come game designer che come narrative. Ma il relatore, per
esempio, sta creando in questo momento un \textbf{card game roguelike
narrative driven, cartaceo}: lì serve un narrative designer, perché un
roguelike ha un \emph{branch} narrativo. E anche in un gioco da tavolo non
ramificato c'è comunque un aspetto di \textbf{user experience} (che cosa il
giocatore legge nel gioco e che cosa capisce) e c'è un \textbf{regolamento}:
``voi non avete idea del ginepraio che è scrivere il regolamento di un gioco
da tavolo''. Di recente il relatore ha fatto il playtest, reale, di un gioco
da tavolo di un'azienda italiana che verrà pubblicato in italiano (``non posso
dire quale''): alla fine ci si è fermati a dare cinquanta minuti di feedback,
``e credo che 45 minuti fossero solo sul regolamento''.

\subsection{L'IA cambierà lo sviluppo? Quali figure rischiano?}

Elena chiede come l'evoluzione dell'IA cambierà lo sviluppo dei videogiochi
nel lungo periodo e quali figure rischiano di diventare superflue.
``Domanda molto, molto, molto complicata.'' Il relatore risponde con ciò che
pensa lui.
\begin{itemize}
  \item \textbf{Primo assunto}: ciò che chiamano IA non è IA, ``è un
        \emph{large language model}''. Un'IA sarebbe Jarvis; questo è uno
        strumento con una probabilità di sbagliare, ``per quanto bassa'', anzi
        ``enorme anche''.
  \item Secondo lui \textbf{dipenderà tutto dall'anno prossimo e dal 2028},
        perché oggi l'IA viene ``spinta in bocca'' dalle aziende, anche a
        prezzo scontato: quanto costerà in futuro? Quale sarà l'impatto?
  \item La sua esperienza: quest'estate ha programmato per una settimana
        volutamente \textbf{solo con l'IA}, ``e sono uscito distrutto''. Amici
        americani obbligati a usarla in azienda gli scrivono che ``nessuno sa
        più fare niente, nessuno sa più dove andare; quando ci sveglieremo da
        questo incubo sarà troppo tardi''.
  \item Il suo assunto: per ora, chi usa l'IA per fare videogiochi o per
        scrivere romanzi è chi \textbf{non vuole davvero fare quella cosa}.
        ``Se tu vuoi fare quella cosa, la vuoi fare tu.'' L'IA può dare una
        mano: ``non sono un luddista''. Il relatore ha Claude, lo paga 20 euro
        al mese, ``non di più'', e si limita nell'uso. Non lo usa mai per
        scrivere; e quando lo usa per il codice, è per cose che sa già fare,
        ``perché so quando Claude mi dice delle stronzate, e non sono poche le
        volte''.
  \item L'IA, secondo lui, cambierà \textbf{i titoli dei ruoli}, e lo sta già
        facendo: si trovano offerte con titoli strani come \emph{assistant AI
        designer}, ``che è un modo per dire: abbiamo licenziato il game
        designer, lo stiamo riassumendo perché l'IA che abbiamo comprato a 100
        euro al mese non fa esattamente quello che volevamo, perché non c'è
        nessuno che capisce quello che l'IA sta dicendo''.
  \item Più dell'IA in sé, conterà l'\textbf{aspetto economico}: non se la
        bolla esploderà, ma come verrà risolta sul piano degli investimenti,
        ``un ginepraio allucinante'', con un rischio enorme di problemi a
        livello globale. ``Non voglio intristirvi.''
  \item Per ora non vede l'IA usata per sostituire troppe persone: molte
        aziende \textbf{hanno riassunto}; molte delle aziende con cui
        collabora hanno nei contratti \textbf{clausole anti-IA}. Dipende
        dall'azienda: Microsoft obbliga a usarla, ma lui dubita che tutti i
        suoi dipendenti la usino per tutto ciò che Microsoft vorrebbe,
        ``altrimenti non verrebbe fatto niente''. Molti software fatti dall'IA
        e dichiarati come tali ``fanno acqua da tutti i pori'': ``non so
        quanti bug ho avuto in Adobe adesso che viene programmato con l'IA''.
\end{itemize}

Elena aggiunge una preoccupazione: forse alle aziende interessa più
velocizzare i processi che avere qualcuno che sa fare benissimo il proprio
lavoro in tempi più lunghi. È un problema che le aziende stanno notando,
risponde il relatore: \textbf{andare veloci è davvero performance?} Se si va
veloci ma il modello 3D va rifatto otto volte, il tempo risparmiato è solo
stato spostato. ``Questa mania del produrre, produrre, essere efficienti,
essere efficienti\ldots\ non è che ripaghi molto.'' La cosa che sta cambiando
davvero, e che un po' lo spaventa, è che \textbf{le posizioni chiederanno di
più}: se ChatGPT può fare un modello 3D di base, ci si aspetterà che un 3D
artist sappia fare qualcosa di più. Molto dipenderà dal business, e c'è un
lavoro enorme da fare sui \textbf{manager}, che hanno ``la tracotanza di dire:
ho visto a casa del mio amico Jeff Bezos questa IA che fa il videogioco
perfetto'', e poi lo si prova ed è ``una merda colossale''. Molti temevano i
romanzi scritti con l'IA pubblicati su Amazon; oggi, dice il relatore, li
legge lo 0,03\% del pubblico di Amazon\footnote{Il dato è del relatore e non è
stato verificato.}. ``La gente se ne accorge.''

``Questo è un anno di transizione.'' Sull'IA il relatore non ha risposte
certe, ``solo una risposta etica'': usarla per tutto, e soprattutto per
sciocchezze, ``è come avere una Ferrari e correre in mezzo alla città
nell'ora di punta: prima o poi qualcuno si farà male''. Bisogna dotarsi di una
\textbf{``patente'' interna}: è molto favorevole a una patente per l'IA, un
corso senza il quale, se la si usa, si passano ``grane forti''. È una
tecnologia potentissima, ``soprattutto a livello di hacking, terribile'': non
perché permetta di fare hacking a chi non sa farlo, ma perché \textbf{l'hacker
che la sa usare è dieci volte più veloce}. ``Mettete le cose in prospettiva,
non fatevi abbindolare dalle informazioni trovate online, dai giornali.''

\subsection{Che cosa pensa di \emph{Control Resonant}?}

``Gioco della Madonna.'' Da un punto di vista di design critico (``perché io
so essere critico''): un piacere per gli occhi, un level design fenomenale,
un gameplay enorme. Il combattimento è carino, ``forse un po' troppo caotico
fin dall'inizio'': avrebbe preferito un \emph{onboarding} un po' più lento,
perché rispetto al primo \emph{Control} il combattimento è \textbf{in mischia}
e non più a distanza, si va subito in combattimento e si creano ``arene di
caos'' che in parte sono volute, ma rendono poco leggibili gli scontri. E in
un gioco in cui ci si deve costruire la \emph{build}, saper leggere il
combattimento è molto importante. A livello narrativo non gli è piaciuto che
si debba \textbf{tornare dagli NPC per far avanzare la loro storia} (tranne
alcuni): se non si torna spesso da un personaggio, poi ci si devono fare
tutti insieme i dialoghi arretrati, e questo crea dissonanza. Avrebbe
preferito che quei dialoghi si perdessero, ``perché crea un'esperienza
narrativa più reale''.

\subsection{Libri per chi si avvicina al narrative}

Elena (che è di 3D art) chiede libri o video per chi si avvicina al narrative
per la prima volta, per esempio sviluppando un piccolo gioco in team. Agli
studenti di narrative il relatore direbbe di aspettare le lezioni; agli altri
consiglia:
\begin{itemize}
  \item \emph{Narrative Design for Mobile and Live Games}, soprattutto le
        prime cinquanta pagine, ``perché ti spiega esattamente bene che cos'è
        il narrative design moderno'';
  \item un libro di world building ``per scrittori'', ma interessante perché
        basato su un'indagine reale, che dà strumenti tecnici per il world
        building\footnote{Nella trascrizione: ``world building for fantasy,
        fonts and author''. Il titolo esatto (forse un \emph{Worldbuilding for
        Fantasy Fans and Authors}) e l'autore non sono ricostruibili con
        certezza.};
  \item \emph{Collaborative Worldbuilding for Video Games} di Kaitlin
        Tremblay\footnote{La trascrizione dice ``Kathleen Tremblay''.},
        ``un po' superficialotto, però come introduzione è molto buono'', utile
        anche per un 3D artist;
  \item per chi è interessato al game design, \emph{Game Design Tools} di
        Diego Ricchiuti, direttore di game design di AIV\footnote{La
        trascrizione dice ``Diego Ricciuti''. Il libro è \emph{Game Design
        Tools: Cognitive, Psychological, and Practical Approaches}.}, per
        capire esattamente che cosa fa un game designer, ``soprattutto perché
        sono le stesse cose che insegniamo in AIV''.
\end{itemize}

\subsection{La narrativa senza parole: ambiente, suono, musica}

Uno studente osserva che, come si è detto, la narrativa non è solo scrittura
ma anche \emph{ambiance}: ci sono momenti senza nessun dialogo che creano un
sentimento, soprattutto grazie alla musica. ``Assolutamente.'' Un tempo in AIV
c'era anche un corso di musica, ``però era troppo costoso da mantenere''. È
fondamentale lavorare con un buon \textbf{compositore} o un buon \textbf{sound
engineer}, perché c'è tutto l'aspetto dei suoni; e nei videogiochi ``la cosa
fichissima'' è che \textbf{i suoni si possono triggerare}: si può mettere un
suono esattamente in un punto. ``Voi non sapete che bellezza è fare horror
nei videogiochi'': si mette un trigger preciso e il giocatore si chiede ``ma
quella cosa era scriptata oppure no?''.

E l'esempio supremo: \emph{P.T.}, la demo mai diventata gioco di
\emph{Silent Hills}. Come hanno creato quel senso di terrore? Per tutto il
gioco, mentre si cammina, \textbf{cinque metri dietro il giocatore c'è il
mostro che lo segue}, con i suoni e i trigger attivi; l'unica cosa disattivata
sono i \emph{collider}, così il mostro può passare dentro i muri. ``Tutte le
volte che vi girate perché sembra di avere qualcuno dietro, avete davvero
qualcuno dietro.'' ``Questo è un game design elegantissimo, il picco del game
design'', ed è il motivo per cui il relatore ama Hideo Kojima: non tanto per
la narrativa, che a lui non piace molto, ma per queste ``chicche'' (come
entrare nel menu e sbloccare cose dal menu).

\subsection{Costruire un personaggio nel tempo}

Sara riprende l'esempio del boss: se un personaggio parte bene e comincia a
creare \emph{pathos}, è più facile proseguire su quella strada, o è più facile
perdersi man mano che si mette carne al fuoco? ``Bellissima domanda, molto
profonda.'' Ci si può perdere in due modi, che si vedranno nel corso.
\begin{enumerate}
  \item \textbf{La perdita produttiva}, come la stagione 8 di \emph{Game of
        Thrones}: se non si è seminato che un personaggio è folle, ``non puoi
        fare retcon di una cosa che non esiste''. Ma è soprattutto una perdita
        in fase di produzione, perché ``noi siamo entertainment in una gabbia
        produttiva''. Esempio: tutti i dialoghi sono stati doppiati da
        un'attrice famosa; ci sono due battute da rifare, si può richiamarla?
        No, perché è già stata pagata e ha l'agenda bloccata per i prossimi due
        anni: ``e mo' come famo?''. Bisogna trovare una soluzione.
  \item \textbf{La perdita del personaggio}: se non lo si costruisce e non lo
        si evolve nel tempo, se si pensa che resti sempre lo stesso, e
        soprattutto se si dà per scontato di poter \textbf{aggiungere senza mai
        togliere}, si rischia un personaggio che diventa una \textbf{Mary
        Sue}, ``termine tecnico'' per un personaggio con una \emph{plot armor}
        incredibile, che sa fare tutto e non morirà mai (il corrispettivo
        maschile è \textbf{Gary Stu}\footnote{La trascrizione dice ``Mary
        Sioux'', ``Gary Sioux'' e ``Garrison''.}); oppure un personaggio che
        \textbf{si perde}, che non si riconosce più e che non dice più niente.
\end{enumerate}
Per il relatore è ciò che è successo a \textbf{Lara Croft} in \emph{Tomb
Raider}. Lo racconta come gliel'aveva spiegato una compagna delle superiori
parlando di due ragazzi che si erano lasciati: non avevano litigato, ma negli
anni il rapporto si era raffreddato e ``non c'era più significato''. Così il
rapporto fra il pubblico e il personaggio: con Lara Croft hanno fatto un lavoro
di decostruzione e messo tanta carne al fuoco, e lentamente il personaggio non
dice più niente. Un po' lo stesso problema di \emph{Uncharted}; ed è il motivo
per cui, secondo lui, ha senso che in \emph{The Witcher 4} la protagonista sia
\textbf{Ciri}: Geralt ``ormai ha detto tutto''; o lo si reinventa, o si fa un
reboot, ma si rompe la bellezza di un personaggio costruito magistralmente in
\emph{The Witcher 3}. ``Dopo tutto quello che abbiamo dovuto soffrire con
Netflix'', CD Projekt Red ha amore per quel personaggio: ``non roviniamolo,
piuttosto abbracciamo il cambiamento''.

\subsection{L'IA per l'arte, e un'IA etica}

Uno studente propone che l'IA si usi solo dove non si tratta di comporre arte,
anche perché per ora ``rende tutto veramente finto''. ``Assolutamente
corretto.'' Per il relatore l'uso etico dell'IA è duplice.
\begin{enumerate}
  \item \textbf{Un'IA locale, non globale.} Esiste un'IA chiamata Colibrì,
        fatta da due italiani, che si può installare sul proprio portatile:
        ogni azienda dovrebbe avere il suo database, la sua mappa semantica,
        il suo modello, addestrato solo sulle proprie conoscenze\footnote{Il
        nome e le caratteristiche del progetto sono quelli riferiti dal
        relatore e non sono stati verificati.}. Il grosso problema è che l'IA
        che usiamo ``ha rubato la proprietà intellettuale di centinaia di
        migliaia di persone'', e le aziende non ne stanno ancora pagando le
        conseguenze; e chi lavorerà in questo settore firmerà contratti in cui
        cede la proprietà intellettuale all'azienda, per cui la questione lo
        riguarda direttamente. L'IA è però così potente che non si può
        ignorarla: è comoda, \textbf{se chi la usa sa dove fermarsi e sa come
        funziona}; bisogna impararla. Chi ne ha paura vada a vedersi i video di
        \textbf{Andrej Karpathy}\footnote{La trascrizione dice ``Antoni
        Carpati''. Il relatore lo descrive come ``ex OpenAI, che è andato adesso
        ad Anthropic''; il capitolo non ha modo di verificare l'informazione.
        Karpathy ha pubblicato su YouTube lezioni introduttive gratuite sui
        modelli linguistici di grandi dimensioni.}, secondo il relatore
        ``l'unica persona oggettiva nel mondo dell'IA'', che spiega gratis su
        YouTube come funziona un large language model: ``non avete scuse''. Un
        po' tranquillizza, perché fa capire anche i limiti della tecnologia.
  \item \textbf{Processi ``alla Jarvis''.} Il relatore sarebbe
        ``fenomenalmente contento'' di avere un Jarvis come Iron Man: fare il
        caffè e intanto dire al computer ``aprimi Adobe, aprimi il file di
        ieri''. ``Quello che mi serve velocizzare non è il processo, è la
        lavatrice, ragazzi, la fottuta lavatrice'', separare i colori dai
        bianchi; ``e non sempre: a volte voglio portare i miei panni alla
        lavanderia sotto casa''.
\end{enumerate}
Bisogna trovare modi di usare l'IA che non facciano danni, perché ``la cosa
peggiore'' è che ci sono IA che fanno danni enormi e le persone non se ne
accorgono, ``hanno il salame sugli occhi''. E l'IA ha rotto il primo principio
della user experience, \textbf{empatizzare con gli altri}: il relatore odia
che Google abbia messo l'\emph{AI Mode} in alto a sinistra, dove c'era la
ricerca, ``perché io quando vado su Google voglio la ricerca, e io sono
l'utente, e l'utente è il tuo cliente''; non si può imporre una funzione solo
perché ci si è investito miliardi.

Ci sono anche usi sensati: strumenti venduti a migliaia di euro che ci si può
fare in casa. Il relatore si fa molti \emph{plug-in} di Unreal, magari con dei
programmatori; e, siccome fa siti web e programma un po' in JavaScript, si fa
aiutare da Claude per costruirsi plug-in per i fogli di calcolo di Google: è
comodo, perché non deve disturbare il programmatore, che magari sta facendo
cose più importanti. Ma non va usata ``come un tappeto, come un'erba che
mettiamo ovunque'': servono limiti, altrimenti ci si ritroverà con sempre più
``merda piatta'', come le locandine delle sagre di paese tutte fatte con la
stessa IA, nello stesso stile. E ci sono grafiche, amiche del relatore, a cui i
clienti portano una locandina fatta con l'IA da ``fixare'': ``non andava bene,
e adesso non va più bene''.

\subsection{Varie: gli orari, i giochi, i remake}

``Abbiamo ancora dieci minuti.'' Il relatore ricorda gli appuntamenti: il
pomeriggio, dalle 14:30 alle 16, l'incontro su mindset e best practice, solo
per chi si è iscritto entro il 31 luglio con Dev Promo o Gamer Promo, ``un'ora
e mezza su metodi di studio, cose avanzate di neurocognitività''
(capitolo~\ref{p0:cap:mindset}); il giovedì alle 10, per tutti, la
comprensione del testo (capitolo~\ref{pt:cap:comprensione}), che resterà registrata per chi lavora. Se non si
trova il link dell'incontro pomeridiano, lo si chiede alla segreteria (è ``un
evento Growth''). (Aveva detto alle 14, poi sposta alle 14:30: ``perché alle
14 finisco di pranzare, magari tardi''.)

Fra le altre domande:
\begin{itemize}
  \item \emph{1666 Amsterdam}: il relatore non l'ha ancora provato, ``come
        idea mi sembra interessante''; è in accesso anticipato, ``io
        aspetterò''. L'unico gioco che ha giocato in early access è stato
        \emph{Hades II}, che ha finito in early access, ``perché è un gioco
        della Madonna, anche quello''.
  \item \textbf{Remake e remaster che toccano la narrativa di un vecchio
        gioco}: per lui è anzitutto un problema di terminologia, ``potato,
        potahto'': lo chiamerebbe \emph{remake}, non \emph{remaster}. Ci sta,
        magari non stravolgendo del tutto ma aggiustando qualcosa, perché negli
        anni le cose possono cambiare; stravolgere del tutto, no.
\end{itemize}

\subsection{Perché gli indie stanno battendo i tripla A?}

Sara chiede: gli indie, ``a livello di anima'', non solo di narrativa, stanno
superando i tripla A; il relatore è d'accordo, e perché succede? ``Non sono solo
d'accordo, ma so perché sta succedendo'', perché parla con i colleghi che
lavorano in quelle aziende: \textbf{è un problema di management}. ``Abbiamo un
management in crisi, che non sa dare risposte'', CEO e direttori che non sanno
bene che cosa fare; ``pensate a PlayStation, che ogni settimana sembra
spararsi su un piede diverso, e sembra avere otto piedi''.

Fra il 2020 e il 2021 c'è stato un grande aumento di assunzioni, perché
``dannate McKinsey e dannate aziende di consulenza americane'' hanno detto che
l'aumento di fatturato dell'industria durante il COVID sarebbe durato per
anni. In quel momento si sono assunti manager che non venivano dalla game
industry: ``il primo momento in cui, anziché assumere dal basso, si è assunto
dal parallelo''. ``\emph{I kid you not}'': il relatore si è trovato a fare
consulenze con direttori creativi che avevano otto anni di esperienza nella
gestione di hotel. Il cliente gli diceva ``vabbè, ma hotel a cinque stelle'';
``sì, ma fanno hotel, non videogame''. Queste persone (``io non ne faccio una
colpa'') vanno in crisi, sono pagate ``ziglioni'' e danno indicazioni come
possono, a volte contraddittorie. Se dietro c'è un capo che ogni settimana
vede una cosa diversa in TV e dice ``questa settimana cambiamo gioco e andiamo
in quella direzione, no, adesso andiamo nell'altra'', o se c'è una
\textbf{yes culture}, i tripla A ``fanno buchi''. È il problema di
\emph{Concord}, ma anche di \emph{Marathon}: ``noi siamo Bungie, non
preoccupatevi dei playtest negativi, il gioco verrà una figata'', quando gli
sviluppatori sapevano che non sarebbe venuto bene.

È logico che gli indie battano i tripla A. La risposta è una sola:
\textbf{``lasciami fare il mio lavoro''}. ``Mi paghi per una professionalità,
lasciami fare il mio lavoro.'' Invece ci sono troppe persone, ``soprattutto i
capi'', che vogliono fare il lavoro degli altri e metterci il naso. Negli
indie funziona meglio perché si è in pochi: invece del meeting del meeting del
meeting, ci si guarda negli occhi due minuti, ci si accorda e si risolve il
problema.

Gabriele concorda: alcune grandi aziende sembrano ferme ai giochi ``patatina'',
gli open world alla \emph{Far Cry 5}, con una mappa apertissima e punti a caso
sulla mappa, mentre l'indie deve dare il tutto per tutto, e nascono ``figate''
come \emph{Hollow Knight} e altri giochi più piccoli\footnote{La trascrizione
riporta anche ``Caped'' e ``coltop dell'erba al ball'': forse \emph{Cult of the
Lamb} e \emph{Balatro}, ma non è ricostruibile con certezza.}. Il vantaggio
dell'indie, aggiunge il relatore, è anche che \textbf{non si può fare un MMO}.
Quante aziende sono state assunte e poi chiuse per fare un MMO? Cita Bluepoint,
lo studio del remake di \emph{Demon's Souls}, che secondo lui è stato chiuso
dopo aver lavorato a un MMO\footnote{Il relatore riferisce che lo studio è
stato chiuso e collega la chiusura al progetto di un MMO; sono notizie del
periodo, riportate qui come le dice e non verificate.}. ``Eh, però ci hanno
toppato loro, su un MMO'': ``certo che vi hanno toppato un MMO: voi chiedevate
solo MMO, davate soldi solo per gli MMO''. E chi ha dato il \emph{green light}
a un MMO di \emph{Horizon} con quello stile? ``Probabilmente un cinquantenne
maschio caucasico con figli di 13 anni che ha visto che \emph{Fortnite} va
forte'': ``decisioni di pancia''.

\subsection{La fine dei dischi fisici?}

Ultima domanda: che cosa pensa del polverone sulla questione del digitale
(Sony, nel 2028, smetterebbe di produrre i giochi su disco)? ``Serviva un
marketer, a Sony'', ``servirebbero dei lavoratori che fanno il loro lavoro, nei
rami alti e non nei bassi'': è logico che una notizia come ``nel 2028 non
facciamo più il disco'' sollevi un polverone, in un momento in cui le aziende
fanno di tutto per far passare l'acquisto digitale ``come se fosse un noleggio,
che è abominevole''. Perché lo fanno? Perché durante il COVID sono state
consigliate male: ``fossi in loro, andrei a chiedere i danni alle società di
consulenza'' che dicevano che il futuro era il digitale e che le persone
sarebbero rimaste sempre in casa.

Uno studente cita Satoru Iwata: ``sul biglietto da visita era un presidente,
nella testa uno sviluppatore, nel cuore un videogiocatore''. ``Esatto.'' Il
relatore aggiunge che il fatto che il presidente di Sony, il giorno dopo la
dichiarazione, abbia venduto il 49\% delle sue quote ``non ispira molta
fiducia''\footnote{Le notizie su Sony (la fine dei dischi nel 2028, la vendita
di quote) sono quelle riferite dal relatore nell'autunno 2026 e non sono state
verificate.}. Secondo lui PlayStation sta facendo così perché ha capito, ``dati
alla mano'', che la PlayStation 6 non arriverà presto e che, quando uscirà, non
farà i soldi sperati; e quindi deve \textbf{tagliare} (``\emph{cut corners}'').
Chi prende queste decisioni si difenderà dicendo ``abbiamo fatto quello che
potevamo'', davanti a numeri che impongono tagli. È come quando un governo
lascia salire quello dello schieramento opposto per fare i tagli, e poi si
prende il merito di non averli fatti, ``e viceversa'' (``non voglio fare torto
a nessuno''; ``in UK è successo così, per non parlare della politica italiana,
in cui non voglio entrare'').

\subsection{I saluti}

``Abbiamo fatto le 11:31.'' Appuntamento al pomeriggio alle 14:30 e al giovedì
mattina; per i narrative, la prima lezione sarà \textbf{venerdì 9 ottobre},
in presenza a Milano per chi ha preso l'Experience, oppure da remoto, dalle 10
alle 13. ``Grazie a tutti, ciao, a dopo.''
